\section{Design Process or Methodology Overview}

In this phase, we focused on transforming the high-level idea of Civic.ai into a concrete, testable system design. Previously, the work was framed as a citizen-centric AI-powered government service assistant. But now, the design has been refined into an evidence-based Retrieval-Augmented Generation (RAG) system for Bangladeshi government service question answering. We have limited the current implementation scope to the birth registration domain only so that the design can be evaluated carefully before adding additional domains in future such as passport, BRTA or other civic services.

The central design requirement of our system is factual correctness. We expect that a government-service assistant should not respond like a general creative chatbot. Rather, it should retrieve relevant evidence from the trusted documents, generate answers from that evidence, give source information and restrict any unsupported claim. For this reason, RAG based design is used in our Civic.ai system, where a retriever first selects relevant document chunks and a generator then produces an answer grounded in the retrieved evidence \cite{lewis2020rag}.

The design process follows five major stages:

\begin{enumerate}
    \item Collection of raw government-service documents from birth registration sources.
    \item Manual and semi-manual cleaning to remove web noise, duplicate navigation text, OCR artifacts, and formatting inconsistencies.
    \item Document aware chunking for different types such as FAQ, rule, fee-table, correction process and application process documents.
    \item Indexing with both dense multilingual embeddings and sparse lexical retrieval.
    \item Retrieval, fusion, reranking, answer generation and evaluation using curated Bangla scenario-based questions.
\end{enumerate}

\begin{figure}[h]
    \centering
    \includegraphics[width=\textwidth]{updated_p2_architecture_diagram.png}
    \caption{Proposed CivicRAG architecture for the birth registration domain. The offline indexing pipeline is separated from the online question-answering pipeline so that retrieval quality and generation quality can be evaluated independently.}
    \label{fig:civicrag_architecture_updated}
\end{figure}

Figure \ref{fig:civicrag_architecture_updated} shows the overall architecture of our current system. It is divided into two pipelines. The offline pipeline at first, processes raw government documents, cleans and chunks the data. Then, it creates two text versions per chunk, embeds the searchable text and stores it in a vector database. The online pipeline accepts a user query, detects language and civic intent, retrieves candidate chunks using dense and sparse retrieval. After that, it fuses ranked results, reranks evidence and finally generates a grounded response using either controlled answer paths or a local LLM.

\section{Preliminary Design Specification}

We considered and simulated several design solutions. The purpose of this comparison was not just to choose one final system but, additionally to justify why the selected design is appropriate for a factual Bangla government-service assistant.

\subsection{Design Solutions}

\begin{table}[h]
\centering
\caption{Design solutions considered for Civic.ai.}
\begin{tabular}{|p{3.2cm}|p{5.4cm}|p{5.2cm}|}
\hline
\textbf{Design Alternative} & \textbf{Description} & \textbf{Reason for Evaluation} \\
\hline
Dense-only RAG & Uses multilingual dense embeddings and vector similarity search. & Strong candidate for paraphrased Bangla queries where exact words may differ from source documents. \\
\hline
BM25-only RAG & Uses sparse lexical retrieval based on term matching. & Useful for exact government terms, legal rule numbers, fee terms and service names \cite{robertson2009bm25}. \\
\hline
Hybrid RAG & Combines dense retrieval and BM25 using rank fusion. & Provides a more robust architecture for future multi-domain expansion where both semantic and exact matching may matter. \\
\hline
CivicRAG & Adds metadata, source URLs, language and intent detection, reranking, safe answer paths and local LLM generation. & Selected design because it supports explainability, factual grounding and controlled answers for high-risk civic service questions. \\
\hline
\end{tabular}
\end{table}

The dense retrieval component is motivated by prior work on dense passage retrieval for open domain question answering \cite{karpukhin2020dpr} and sentence-level semantic representation learning \cite{reimers2019sbert}. In our current implementation, we used the BAAI/bge-m3 embedding model as it is designed for multilingual, multi-functionality and multi-granularity retrieval \cite{chen2024bgem3}. This is important for our Civic.ai system because our knowledge base is Bangla-heavy while user queries may appear in Bangla, English or code-mixed form.

The sparse retrieval component uses BM25, a strong probabilistic lexical retrieval baseline \cite{robertson2009bm25}. BM25 is specifically useful when a query contains exact service words, fee-related words, rule references, or official terminology. The simplified BM25 scoring function can be represented as:

\begin{equation}
BM25(q,d)=\sum_{t \in q} IDF(t)\cdot
\frac{f(t,d)(k_1+1)}
{f(t,d)+k_1(1-b+b\cdot\frac{|d|}{avgdl})}
\end{equation}

Here, $q$ is the query, $d$ is a document chunk, $f(t,d)$ is the frequency of term $t$ in chunk $d$, $|d|$ is the chunk length and $avgdl$ is the average chunk length. The parameter $k_1$ controls term-frequency and $b$ controls saturation and length normalization \cite{robertson2009bm25}.

To combine dense and sparse retrieval results, the system uses Reciprocal Rank Fusion (RRF). RRF combines multiple ranked lists without requiring the raw scores of different retrievers to be directly comparable \cite{cormack2009rrf}. In the current system, a weighted form of RRF is used:

\begin{equation}
RRF(d)=\sum_{r \in R}\frac{w_r}{k+rank_r(d)}
\end{equation}

where $d$ is a candidate chunk, $R$ is the set of retrievers, $w_r$ is the retriever weight, $k$ is the RRF constant, and $rank_r(d)$ is the rank of the chunk $d$ under retriever $r$. The original RRF formulation is unweighted, but the weighted version is useful for Phase 2 sensitivity analysis because it allows the dense and BM25 contributions to be tuned.

After going through fusion, a reranking stage is used to improve the order of candidate chunks. Cross encoder reranking is a common approach for passage reranking because it jointly scores the query and candidate passage \cite{nogueira2019passagebert}. The current Civic.ai implementation supports cross-encoder reranking when the local model is available and uses deterministic lexical/domain-aware fallback reranking for local reproducibility.

\section{Data Collection}

The current dataset focuses on Bangladeshi birth registration services. The sources include government service web pages, frequently asked questions, legal rules, fee information, application process documents, correction guidance and scanned or converted documents related to birth registration.

Four major raw data types were considered:

\begin{itemize}
    \item Bijoy-encoded PDF files.
    \item Converted DOCX files.
    \item Unicode HTML pages from government web portals.
    \item Scanned PDF or OCR-affected documents.
\end{itemize}

The current Phase 2 implementation primarily indexes manually reviewed and cleaned Markdown/JSON data. Raw scanned and converted files are retained for auditability, but only cleaned and structured content is used for the active vector database. This decision was made to reduce garbage-in-garbage-out errors in retrieval.

\subsection{Data Cleaning}

Civic service documents are often noisy and mixed or diverse. HTML files often contain menus, accessibility controls, footers, lists of unrelated services and navigation text. PDF and DOCX files sometimes contain broken formatting, repeated headings, encoding issues and OCR noise. Sometimes they include diagrams as well. Data cleaning is considered and treated as one of the core design steps rather than something minor since the quality of RAG heavily depends on the quality of retrieved data.

The data cleaning process includes:

\begin{itemize}
    \item Removing menus, footer blocks, accessibility controls and unrelated service lists.
    \item Removing gibberish texts and symbols and irrelevant text.
    \item Removing and converting all the non-Unicode/Bijoy texts where possible.
    \item Turning FAQ entries into individual question-answer units for creating safe answering paths.
    \item Preserving legal and procedural wording where exact phrasing is important.
    \item Eliminating OCR fragments, empty headings and duplicate content.
    \item Preserving official source URLs as metadata for explainability.
\end{itemize}

The goal of cleaning data is to prevent irrelevant and unnecessary texts from being embedded and retrieved. If noisy navigation or meaningless text is stored in the vector database the retriever may return chunks that match common words but do not answer the citizen's question or include gibberish texts in the response.

\begin{figure}[h]
    \centering
    \includegraphics[width=0.88\textwidth]{updated_p2_data_cleaning_diagram.png}
    \caption{Data cleaning and preprocessing pipeline for proposed CivicRAG.}
    \label{fig:data_cleaning_pipeline_updated}
\end{figure}

\subsection{Data Transformation}

The data is transformed into retrieval friendly useful chunks after the cleaning process. Since there are various kinds of document types and one single chunking method is not suitable for that, we are using document-type-aware chunking.

\begin{table}[h]
\centering
\caption{Chunking strategy by document type.}
\begin{tabular}{|p{3.5cm}|p{5.1cm}|p{5cm}|}
\hline
\textbf{Document Type} & \textbf{Chunking Strategy} & \textbf{Reason} \\
\hline
FAQ JSON & Question-answer pair wise JSON chunk. & FAQ answers the most accurate answers and it is better to keep them intact. \\
\hline
Heading based Markdown & Section and subsection based chunking. & Preserves procedural structure and improves evidence interpretability. \\
\hline
Fee tables & Row wise chunks with fee metadata. & Fee questions require precise row level retrieval. \\
\hline
Legal/rule documents & Rule or paragraph group based chunks. & Legal evidence should be separated from unrelated sections. \\
\hline
Correction/application process documents & Step or condition based chunks. & Queries often include specific scenarios or require procedural steps. \\
\hline
\end{tabular}
\end{table}

Each chunk contains two separate text fields:

\begin{itemize}
    \item \textbf{retrieval\_text:} used for searching. It may include aliases, normalized terms and key words.
    \item \textbf{answer\_evidence:} used for generation. It contains original source content and is passed to the LLM for grounding.
\end{itemize}

This separation is required because aliases can elevate retrieval quality but the generator should not be grounded on artificial aliases. Therefore, the LLM receives original source content rather than expanded retrieval text.

\subsection{Data Integration}

All domain specific resources are organized under one domain. The current domain is:

\begin{center}
\texttt{domains/birth\_death\_registration/}
\end{center}

This folder holds cleaned Markdown/JSON files, generated chunks, ChromaDB vector storage, configuration files and evaluation datasets including raw documents. This design will help future multi domain expansion without having to create separate repositories. Future domains such as passport or BRTA as discussed can be added as new domain folders with their own data and configuration files.

\subsection{Data Reduction}

In this project, data reduction does not actually mean numerical dimensionality reduction. Instead, it means reducing irrelevant document noise and gibberish texts before indexing and reducing retrieved candidate chunks before generation. The important reduction decisions are:

\begin{itemize}
    \item Eliminating duplicate and irrelevant government website boilerplate.
    \item Removing duplicate FAQ content where the same question/answer appears in multiple sources.
    \item Splitting long documents into focused evidence units.
    \item Using fine tuned top-k retrieval to avoid sending too much context to the generator.
    \item Applying reranking and intent based boosting to prioritize the most relevant evidence instead of blindly sending the retrieved evidence.
\end{itemize}

Data Reduction is necessary because long or noisy retrieved data can confuse LLM generation. In a government service system, the answer should be precise and provide enough information but controlled because the actually answer should not be influenced or overshadowed by the unrelated chunks.

\subsection{Summary of Preprocessed Data}

The final preprocessed data contains Bangla heavy birth registration chunks. The major chunk categories are FAQ, legal rules, fee rows, correction steps, application steps, portal guidance and general service information. Each chunk holds metadata such as source ID, document type, section title, category and source URL.

For Phase 2 evaluation, a 58-question Bangla paraphrase evaluation set was created from 19 grouped birth certificate related problems. These questions are designed to evaluate whether the system can retrieve the same expected evidence when a citizen asks the same problem in different wording and even if it does how well it can do it.

\section{Implementation of Selected Design}

The selected design is CivicRAG, a modular local RAG system for government service question answering. The implementation consists of six main modules:

\begin{enumerate}
    \item \textbf{Data processing module:} reads cleaned Markdown and JSON files and converts them into typed chunks.
    \item \textbf{Embedding module:} embeds \texttt{retrieval\_text} using BAAI/bge-m3 \cite{chen2024bgem3}.
    \item \textbf{Indexing module:} stores dense vectors in ChromaDB and builds a BM25 sparse index.
    \item \textbf{Retrieval module:} retrieves candidates using dense search, BM25 search, or hybrid RRF search.
    \item \textbf{Reranking and safety module:} reranks candidates and checks whether a controlled answer path is available.
    \item \textbf{Generation module:} produces a final answer using local LLMs when a controlled answer is not available.
\end{enumerate}

The current local LLM backends are llama3.2, llama3, and qwen2.5:7b. Llama-based models are used as local open-weight baselines \cite{dubey2024llama3}, while Qwen2.5 is included as an alternative multilingual local model \cite{qwen2025qwen25}. The default generation temperature is set to 0.05 since government service queries require precise and consistent answers rather than creative answers. This low temperature eliminates the randomness.

For high-confidence civic intents such as fee questions, lost certificate questions, correction questions and application-process based questions the system can use safe answer paths. These paths reduce unnecessary LLM intervention by producing answers directly from selected evidence patterns. Retrieved evidence is sent to the local LLM with a strict grounding prompt to generate an answer if no safe path is found.

Generation quality is decided to be evaluated using metrics such as answer correctness, answer relevance, faithfulness and qualitative human review. RAG-specific evaluation frameworks such as RAGAS emphasize faithfulness, answer relevancy and context quality for retrieval augmented systems \cite{es2024ragas}.

Therefore, the Phase 2 implementation satisfies the design requirement of evaluating multiple alternative solutions before selecting the final architecture. Dense-only retrieval currently has the highest score on all fields, therefore, is the strongest retrieval baseline on the single domain Bangla evaluation set. Although CivicRAG remains the chosen system design since it is more explainable, supports source grounded responses and is better prepared for future multi-domain civic service expansion.

\section{Preliminary Evaluation}

A curated Bangla evaluation set was created to simulate real citizen questions. The current evaluation set contains 58 paraphrased questions. These questions cover application processes, correction cases, lost certificate cases, fees, legal conditions, online visibility problems and document requirements.

The retrieval alternatives were evaluated using ranking metrics commonly used in information retrieval, including Hit@k/Recall@k, Mean Reciprocal Rank (MRR), normalized Discounted Cumulative Gain (nDCG) and Context Precision \cite{manning2008ir,voorhees1999trec8qa,jarvelin2002ndcg}. The main retrieval results are shown in Table \ref{tab:retrieval_comparison_updated}.

\begin{table}[h]
\centering
\caption{Retrieval variant comparison on the 58-question Bangla evaluation set.}
\label{tab:retrieval_comparison_updated}
\begin{tabular}{|l|c|c|c|c|}
\hline
\textbf{Method} & \textbf{Hit@1} & \textbf{Hit@5} & \textbf{MRR} & \textbf{nDCG@5} \\
\hline
BM25-only & 0.5345 & 0.7586 & 0.6132 & 0.6112 \\
\hline
Dense-only & 0.7931 & 0.9828 & 0.8707 & 0.8693 \\
\hline
Hybrid RRF & 0.7414 & 0.9483 & 0.8193 & 0.8139 \\
\hline
\end{tabular}
\end{table}

The result shows that dense-only retrieval currently performs best for the birth registration evaluation set. This is reasonable because many evaluation questions are slightly modified or paraphrased Bangla questions, where semantic similarity is more important than exact word overlap. However, hybrid retrieval remains part of the selected design because future multi-domain expansion may increase lexical ambiguity. For example, different government services may contain semantically similar instructions but domain-specific terms such as passport, driving license, birth registration, death registration, fee, correction, or application. In such settings, sparse retrieval can help preserve exact domain signals.

The following metrics were used:

\begin{equation}
Recall@k = \frac{\text{Number of queries where a relevant chunk appears in top }k}{\text{Total number of queries}}
\end{equation}

\begin{equation}
MRR = \frac{1}{|Q|}\sum_{i=1}^{|Q|}\frac{1}{rank_i}
\end{equation}

MRR measures how early the first relevant chunk appears in the ranked list \cite{voorhees1999trec8qa,manning2008ir}. For ranking quality with graded relevance, nDCG is used:

\begin{equation}
DCG@k = \sum_{i=1}^{k}\frac{2^{rel_i}-1}{\log_2(i+1)}
\end{equation}

\begin{equation}
nDCG@k = \frac{DCG@k}{IDCG@k}
\end{equation}

where $rel_i$ is the relevance score of the result at rank $i$, and $IDCG@k$ is the ideal DCG at rank $k$ \cite{jarvelin2002ndcg}.

The generation side design was also examined through context size testing. The goal was to determine how many retrieved chunks should be sent to the answer generator. Table \ref{tab:topk_context_updated} shows the trade-off between evidence coverage and context precision.

\begin{table}[h]
\centering
\caption{Generation context top-k coverage.}
\label{tab:topk_context_updated}
\begin{tabular}{|c|c|c|}
\hline
\textbf{k} & \textbf{Expected Source in Context} & \textbf{Context Precision@k} \\
\hline
1 & 0.7414 & 0.7414 \\
\hline
3 & 0.8793 & 0.3621 \\
\hline
5 & 0.9483 & 0.2448 \\
\hline
8 & 0.9828 & 0.1616 \\
\hline
\end{tabular}
\end{table}

The results show that larger context sizes improve the chance of including the expected evidence but also introduce more irrelevant chunks. The current selected value is $k=3$ for generation because it is conservative and reduces noise. However, $k=5$ is a candidate for future testing because it improves expected-source coverage.
